% Append these entries inside the final paper's thebibliography environment.
\bibitem{CAlpoge26} L. Alp\"oge, \emph{Remark}, supplied four-page note, 20 November 2026. \url{https://alpo.ge/remark.pdf}. Supplied \texttt{ALPOGE\_Main.lean}: \texttt{riemannHypothesis}, \texttt{riemannZeta\_eq\_zero\_re}, \texttt{LFunction\_ne\_zero\_of\_half\_lt\_re}, \texttt{LFunction\_eq\_zero\_re}, \texttt{LFunction\_eq\_zero\_re\_of\_isPrimitive}; types in \texttt{ALPOGE\_Statements.lean}. Used as cited premises as requested, without claiming a Lean rebuild.
\bibitem{CKoiran96} P. Koiran, \emph{Hilbert's Nullstellensatz is in the polynomial hierarchy}, J. Complexity 12 (1996), 273--286, \url{https://doi.org/10.1006/jcom.1996.0019}. Exact numbering here is that of the DIMACS long version, Technical Report 96-27: Theorem 2; Section 5, Theorem 9 and Corollary 1. \url{https://perso.ens-lyon.fr/pascal.koiran/Publis/DIMACS.96-27.ps}. Erratum, 8 August 1997: \url{https://perso.ens-lyon.fr/pascal.koiran/Publis/hn.erratum.ps}.
\bibitem{CRojas03} J. M. Rojas, \emph{Dedekind zeta functions and the complexity of Hilbert's Nullstellensatz}, arXiv:math/0301111v2 (2003). Theorems 1 and 2. \url{https://arxiv.org/abs/math/0301111}.
\bibitem{CGCTV} K. D. Mulmuley, \emph{Geometric Complexity Theory V: Efficient algorithms for Noether Normalization}, arXiv:1209.5993v5. Theorems 1.7, 1.9, 2.10, 3.4, 4.1, 4.7, 4.9, 4.10, 5.11, 5.14. \url{https://arxiv.org/abs/1209.5993}. The supplied LaTeX and the retrieved PDF use the numbering cited here.
\bibitem{CMiller76} G. L. Miller, \emph{Riemann's hypothesis and tests for primality}, J. Comput. System Sci. 13 (1976), 300--317. Theorem 2 and Appendix. \url{https://www.cs.cmu.edu/~glmiller/Publications/Papers/Mi76.pdf}.
\bibitem{CLenstra79} H. W. Lenstra, Jr., \emph{Miller's primality test}, Inform. Process. Lett. 8 (1979), 86--88. Theorem 1. \url{https://pub.math.leidenuniv.nl/~lenstrahw/PUBLICATIONS/1979e/art.pdf}.
\bibitem{CBach90} E. Bach, \emph{Explicit bounds for primality testing and related problems}, Math. Comp. 55 (1990), 355--380. Section 6, Theorems 2--3, pp.~369--373; Section 7, Theorem 4. \url{https://doi.org/10.1090/S0025-5718-1990-1023756-8}.
\bibitem{CAKS04} M. Agrawal, N. Kayal and N. Saxena, \emph{PRIMES is in P}, Ann. of Math. 160 (2004), 781--793. Theorems 4.1 and 5.1 (also Theorem 5.3 for the improved bound). \url{https://annals.math.princeton.edu/wp-content/uploads/annals-v160-n2-p12.pdf}.
\bibitem{COAI26} OpenAI, \emph{The Exponential Time Hypothesis is False}, OpenAI Math Release preprint, 17 October 2026, revised 25 October 2026. Theorem 1.1, Corollaries 1.2--1.4, Sections 1.1.1--1.1.3, Appendix C.1, Appendix D, and the supplied files \texttt{quantitative-reduction.tex}, \texttt{quantitative-contraction.tex}, \texttt{quantitative-circuits.tex}. \url{https://github.com/openai/math/blob/main/preprints/The-Exponential-Time-Hypothesis-is-False-October-25-2026/paper.pdf}.
\bibitem{CCGGM26} N. Chukhin, K. Gajulapalli, A. Golovnev and I. Mihajlin, \emph{A Little Derandomization Goes a Long Way: Unconditional Circuit Lower Bounds using Low-Degree PIT}, arXiv:2611.07972v1, 5 November 2026; ECCC TR26-313, published 9 November 2026. Theorems 1, 5 and 6; Sections 5 and 8.2. \url{https://arxiv.org/abs/2611.07972}; \url{https://eccc.weizmann.ac.il/report/2026/313/}.
\bibitem{CLandsberg13} J. M. Landsberg, \emph{Geometric complexity theory: an introduction for geometers}, arXiv:1305.7387v3, 11 December 2013. Section 8.2, equation (22), p.~28. \url{https://arxiv.org/abs/1305.7387}.
